% Related Work
% Organized by subfield, not chronologically

\section{Related Work}
\label{sec:related}

% ============================================================
% 2.1 KV Cache Optimization
% ============================================================
\subsection{KV Cache Optimization}

LLM serving incurs substantial memory overhead from KV caches.
FlashAttention~\cite{dao2022flashattention} addresses this via IO-aware kernel fusion,
reducing memory footprint without changing the KV format.
Speculative decoding~\cite{leviathan2023fast,chen2023accelerating}
reduces latency by predicting multiple tokens ahead,
but requires a draft model to generate candidates.

Our work is orthogonal: we study what transfers across models,
not how to optimize cache memory within a single model.

% ============================================================
% 2.2 Model Compression and Distillation
% ============================================================
\subsection{Model Compression and Distillation}

Model compression reduces size via quantization~\cite{frantar2023gptq,dettmers2024qlora},
pruning, or distillation~\cite{touvron2023llama}.
These methods transform model weights,
while we study transferring KV cache state across existing models.

Knowledge distillation~\cite{hinton2015distilling}
transfers knowledge from teacher to student via soft labels.
Our work shows that KV state can transfer directly,
without the full distillation pipeline.

% ============================================================
% 2.3 Attention Mechanisms
% ============================================================
\subsection{Attention Mechanisms}

The Transformer architecture~\cite{vaswani2017attention}
decomposes attention into keys (addressing) and values (content).
This decomposition is fundamental to our work:
we study whether K and V transfer differently across models.

Prior work treats KV as monolithic.
We show K and V serve different functional roles:
K is model-agnostic routing; V is weight-bound content.

% ============================================================
% 2.4 Model Stitching
% ============================================================
\subsection{Model Stitching}

Model stitching combines layers from different models
to create hybrid architectures~\cite{bansal2020could}.
Recent work shows that attention layers can be stitched
across models with minimal performance loss.

Our heterogeneous reassembly (teacher K + student V)
extends this line:
addressing geometry is composable,
while content representations are not.
This suggests a functional decomposition for model stitching.

% ============================================================
% 2.5 Scaling Laws
% ============================================================
\subsection{Scaling Laws}

Scaling laws~\cite{kaplan2020scaling,hoffmann2022chinchilla}
characterize how performance scales with compute, data, and parameters.
Our cost crossover analysis complements this:
when is it cheaper to transfer state vs.\ transfer weights?
The crossover at $\sim$2050 tokens provides a practical decision rule.

% ============================================================
% 2.6 Summary
% ============================================================
\subsection{Summary}

Our work is distinct from prior efforts in three ways:
\begin{enumerate}
  \item \textbf{Decomposition.} We decompose KV into K and V,
    showing asymmetric transferability.
  \item \textbf{Functional analysis.} We study what transfers functionally,
    not just geometrically (CCA correlation).
  \item \textbf{Cost model.} We provide a practical cost crossover
    for state vs.\ weight transfer.
\end{enumerate}
